\section{Robustness}\label{sec:robustness}
Let $F:\R^n\to\R^m$ be a classifier with $m\ge2$ scores and top-score tie locus
\[
\Sigma=\bigcup_{i<j}(C_i\cap C_j).
\]
Writing $g_{ij}=F_j-F_i$ and $V_{ij}=\mathcal Z(g_{ij})$, we have
\begin{equation}\label{eq:sigma-inclusion}
\Sigma\subseteq\bigcup_{i<j}V_{ij}.
\end{equation}
The inclusion may be strict because $V_{ij}$ also records ties between nonmaximal scores. Our estimates control proximity to this larger union.

Under the regularity assumptions below, each $V_{ij}$ is closed and has no boundary, so $T(V_{ij},\varepsilon)=T^\perp(V_{ij},\varepsilon)$. The normal-tube estimates therefore control the full distance neighbourhoods needed here, even though $\Sigma$ itself may be nonsmooth.

For $x\notin\Sigma$, put
\[
\Delta(x)=\dist(x,\Sigma).
\]
This is the distance to loss of a unique maximal score. Under the pairwise regularity assumptions below, it is also the infimum of perturbation sizes that change the predicted label. Indeed, starting from $x\in\inter C_j$, any path to another label meets $\partial C_j\subseteq\Sigma$, while arbitrarily close to each point of $\partial C_j=\Sigma_j$ there are points outside $C_j$. A segment from $x$ to any point of $\Sigma$ either ends on $\partial C_j$ or meets it earlier. Thus the distance to $\Sigma$ equals the distance to $\partial C_j$, and hence the infimum of distances to a different label. This concerns label stability, not agreement with ground truth.

\begin{definition}\label{def:condition}
The \emph{condition number} of the classifier at $x\notin\Sigma$ is
\[
\mC(x)=\mC_0(x)=\frac{\|x\|}{\Delta(x)}.
\]
More generally, its \emph{local condition number} centred at $p\in\R^n$ is
\[
\mC_p(x)=\frac{\|x-p\|}{\Delta(x)}.
\]
\end{definition}

We use $\dist(x,\varnothing)=\infty$, with the corresponding condition number equal to zero. Under pairwise regularity, $\Sigma$ has Lebesgue measure zero. The condition number may be extended by $\infty$ on $\Sigma\setminus\{p\}$ and defined arbitrarily at $p$; these choices do not affect the probabilities considered below.

For $x\ne p$ off $\Sigma$, the inequality $\mC_p(x)>t$ means that the unique predicted label can be lost under a perturbation of relative size less than $1/t$, with size measured relative to $\|x-p\|$. This parallels the interpretation of distance-based condition numbers in numerical analysis~\cite{burgisser2013condition}.

\begin{remark}[Weak condition number]\label{rem:weak-condition}
For a uniformly random direction $v\in S^{n-1}$, define the first distance to a tie by
\[
\Delta(x;v)=\inf\{\eta>0:x+\eta v\in\Sigma\},
\]
with value $\infty$ if the ray does not meet $\Sigma$. For $0<\delta<1$, the quantile
\[
\mC^{\mathrm{weak}}_\delta(x)=\inf\left\{\eta\ge0:
\prob_v\left\{\frac{\|x\|}{\Delta(x;v)}>\eta\right\}\le\delta\right\}
\]
gives a directional analogue of the weak condition numbers in~\cite{lotz2020weak}. 
\end{remark}

\subsection{Tail bounds on the condition number}
The estimates below apply both to random inputs and to random perturbations of a fixed centre. The local normalisation makes the bounds independent of the sampling radius.

\begin{theorem}[Uniform tail bound]\label{thm:condition-uniform}
Suppose every pairwise difference $g_{ij}:\R^n\to\R$ is Pfaffian of format $(\alpha,\beta,s)$, with $\alpha,\beta\ge1$, and has $0$ as a regular value. If $X$ is uniform in $B(p,\rho)$, then for every $t>0$,
\begin{equation}\label{eq:general-condition-tail}
\prob\{\mC_p(X)>t\}\le
\min\left\{1,\binom m2 C_{\alpha,\beta,s,n}
\left[\left(1+\frac{\alpha+\beta}{t}\right)^n
-\left(1+\frac1t\right)^n\right]\right\},
\end{equation}
where $C_{\alpha,\beta,s,n}$ is defined in~\eqref{eq:pfaffian-tube-constant}. The hypersurfaces $V_{ij}$ may be unbounded.
\end{theorem}

\begin{proof}
On $B(p,\rho)$, the event $\mC_p(X)>t$ implies $\Delta(X)<\rho/t$. By~\eqref{eq:sigma-inclusion}, it is contained, up to a null set, in
\[
\bigcup_{i<j}\{\dist(X,V_{ij})<\rho/t\}.
\]
Apply Theorem~\ref{thm:prob_bound} to each pair with $\varepsilon=\rho/t$, sum the bounds, and use that a probability is at most one.
\end{proof}

A Gaussian distribution is an exact mixture of uniform-ball distributions. We state this independently of the set to which it will be applied.

\begin{lemma}\label{lem:gauss-to-uniform}
Let $X\sim\mathcal N(p,\varsigma^2 I_n)$, with $\varsigma>0$, and let $U_r$ be uniform in $B(p,r)$. For every measurable $A\subseteq\R^n$,
\begin{equation}\label{eq:gauss-to-uniform}
\prob\{X\in A\}
=\frac{\omega_n}{(2\pi)^{n/2}}
\int_0^\infty\prob\{U_{\varsigma r}\in A\}
\,r^{n+1}\econst^{-r^2/2}\,\diff r.
\end{equation}
In particular, any upper bound for $\prob\{U_r\in A\}$ valid for every $r>0$ also bounds $\prob\{X\in A\}$.
\end{lemma}

\begin{proof}
Substitute
\[
\prob\{U_{\varsigma r}\in A\}
=\frac{1}{\omega_n(\varsigma r)^n}
\int_A\mathbf1_{\{\|x-p\|\le\varsigma r\}}\,\diff x
\]
into the right-hand side and apply Tonelli's theorem. The inner integral is
\[
\int_{\|x-p\|/\varsigma}^\infty r\econst^{-r^2/2}\,\diff r
=\exp\!\left(-\frac{\|x-p\|^2}{2\varsigma^2}\right),
\]
giving the Gaussian density on $A$. Taking $A=\R^n$ proves the normalisation.
\end{proof}

\begin{theorem}[Gaussian tail bound]\label{thm:condition-gaussian}
Under the pairwise format and regularity assumptions of Theorem~\ref{thm:condition-uniform}, the bound~\eqref{eq:general-condition-tail} also holds for $X\sim\mathcal N(p,\varsigma^2 I_n)$, for every centre $p$ and every $\varsigma>0$.
\end{theorem}

\begin{proof}
Apply Lemma~\ref{lem:gauss-to-uniform} to the fixed set $A=\{x:\mC_p(x)>t\}$. Theorem~\ref{thm:condition-uniform} bounds its probability in every ball centred at $p$ by the same right-hand side.
\end{proof}

For fixed formats and $t\to\infty$, the expression inside the minimum in~\eqref{eq:general-condition-tail} has expansion
\[
\binom m2 C_{\alpha,\beta,s,n}\,
\frac{n(\alpha+\beta-1)}{t}+O_{\alpha,\beta,s,n,m}(t^{-2}).
\]

\subsection{Neural network classifiers}
Proposition~\ref{prop:nn-format} makes the general estimates explicit for neural networks.

\begin{corollary}\label{cor:nn-condition}
Let $F:\R^n\to\R^m$ have $\ell$ hidden layers and $h$ hidden units, with autonomous activation format $(\alpha,\beta,s)$, $\alpha,\beta,s\ge1$. Assume every pairwise score difference has $0$ as a regular value. For either $X$ uniform in $B(p,\rho)$ or $X\sim\mathcal N(p,\varsigma^2 I_n)$, the bound~\eqref{eq:general-condition-tail} holds with format $(\alpha_F,\beta_F,s_F)$ given by
\[
\begin{cases}
\bigl(\ell(\alpha+\beta-1)-\beta+1,\ \beta,\ sh\bigr),
&\text{for affine outputs},\\
\bigl(\ell(\alpha+\beta-1)+\alpha',\ \beta',\ sh+s'm\bigr),
&\text{for autonomous outputs},
\end{cases}
\]
where in the second case each output component has format $(\alpha',\beta',s')$ with positive parameters.
\end{corollary}

\begin{proof}
All pairwise differences share the common chain in Proposition~\ref{prop:nn-format}. Apply Theorems~\ref{thm:condition-uniform} and~\ref{thm:condition-gaussian}. If softmax follows affine outputs, apply the bounds to the logit differences $\widetilde g_{ij}$: their zero sets and regularity agree with those of the score differences $g_{ij}$, as established in Section~\ref{sec:tubular-neural}.
\end{proof}

\begin{example}[Sigmoid or tanh activation]\label{ex:sigmoid-condition}
For sigmoid or tanh hidden activations and affine outputs, $(\alpha_F,\beta_F,s_F)=(2\ell,1,h)$. Put
\[
B_{\ell,h,n}=1+(n-1)(2\ell-1)+2\ell\min\{n,h\}.
\]
Under pairwise regularity, for either sampling model,
\begin{align*}
\prob\{\mC_p(X)>t\}
&\le\min\left\{1,\binom m2
\frac{2^{h(h-1)/2}}{\ell}B_{\ell,h,n}^{\,h}
\left[\left(1+\frac{2\ell+1}{t}\right)^n
-\left(1+\frac1t\right)^n\right]\right\}.
\end{align*}
Before clipping, its leading term is
$\binom m2\,2^{h(h-1)/2+1}nB_{\ell,h,n}^{\,h}/t$.
\end{example}

For a single sigmoid layer, the refined tube bound gives a polynomial dependence on width for both sampling models, without a containment assumption.

\begin{corollary}\label{cor:nn-condition-single-layer}
Let $F:\R^n\to\R^m$ have one hidden layer of width $w\ge1$, logistic sigmoid activation and affine outputs, optionally followed by softmax. Suppose its first-layer weights $a_k$ have a common denominator $q\in\N_{>0}$, with $L=\max_k\|qa_k\|_\infty\ge1$, and every pairwise logit difference has $0$ as a regular value. For every $p\in\R^n$ and every $t>0$,
\begin{equation}\label{eq:single-layer-condition-bound}
\prob\{\mC_p(X)>t\}
\le\min\left\{1,\binom m2 H(n,L)\frac{w^n}{t}\right\},
\end{equation}
both for $X$ uniform in $B(p,\rho)$ with any $\rho>0$ and for $X\sim\mathcal N(p,\varsigma^2 I_n)$ with any $\varsigma>0$. Here $H(n,L)$ is the explicit constant in~\eqref{eq:single-layer-H}.
\end{corollary}

\begin{proof}
Every pairwise logit difference is of the form
\[
\widetilde g_{ij}(x)=c_{ij}+\sum_{k=1}^w d_{ij,k}\sigma(a_k^{\trans}x+b_k)
\]
with the same first-layer weights and lattice constant, and $V_{ij}=\mathcal Z(g_{ij})=\mathcal Z(\widetilde g_{ij})$. For uniform sampling, apply~\eqref{eq:single-layer-linear-tube} with $\varepsilon=\rho/t$ to each pair, and use the union bound as in Theorem~\ref{thm:condition-uniform}. The result is independent of $\rho$. Lemma~\ref{lem:gauss-to-uniform} therefore gives exactly the same estimate for Gaussian sampling.
\end{proof}

\begin{remark}\label{rem:single-layer-condition-tight}
A more detailed tail profile follows by replacing $u$ with $1/t$ in~\eqref{eq:single-layer-refined-tube} and multiplying by $\binom m2$, for either sampling model. In particular, for $t\ge1$ the resulting upper bound is
\[
\binom m2\left[
\frac{2n((2Lw+1)^n-1)}{t}
+O_{n,L}\!\left(\frac{w^{2n}}{t^2}\right)\right].
\]
The clipped estimate~\eqref{eq:single-layer-condition-bound} is valid for every $t>0$; its nontrivial regime starts at a constant multiple of $w^n$ when $n,L,m$ are fixed.
\end{remark}
